\section{排序、选择和分治的系统讲解}

排序题表面是在排列元素，实际是在制造顺序信息。没有顺序时，两数比较、区间合并、去重、前驱后继都可能需要线性寻找；排序以后，许多关系会变成相邻关系或单调关系。三数之和能够从三层枚举降到固定一数加双指针，是因为排序把剩余两数的和变成可单调调整的量。合并区间能够线性扫描，是因为按起点排序后，所有可能与当前结果尾部相交的区间都会连续出现。排序不是免费的，复杂度通常是 O(n log n)，所以必须说明排序换来了什么查询能力。

比较排序的下界来自信息论：n 个不同元素有 n! 种可能顺序，比较一次最多提供一位左右的信息，因此最坏至少需要数量级为 n log n 的比较。这个结论告诉我们，想突破 n log n，必须利用额外条件，例如值域有限、元素是整数、只需要第 K 个而不是完整顺序、或者允许随机化平均复杂度。计数排序利用小值域，桶排序利用分布，基数排序利用数字位，快速选择利用只关心目标排名。这些算法不是比标准排序更高级，而是使用了更强前提。

快速排序的核心是 partition：选择一个枢轴，把小于它的元素放一边，大于它的元素放另一边。partition 完成后，枢轴位置已经确定，左右两边再分别排序。快速选择只进入目标排名所在的一边，因此平均时间为线性。这里的平均性依赖枢轴足够随机；若每次都选到极端枢轴，最坏会退化到平方。面试中写快速选择必须主动说明随机化或三数取中，不然复杂度承诺不完整。

归并排序和快速排序的思维不同。归并排序先分再合，合并两个有序段时保证稳定性，最坏复杂度稳定为 O(n log n)，但需要额外空间。链表排序常用归并，因为链表合并不需要随机访问，节点重连成本低；数组排序常用 introsort 一类混合策略，在快速排序、堆排序、插入排序之间切换，兼顾平均性能和最坏情况。理解这些取舍后，才能解释为什么 C++ 的 std::sort 不承诺稳定，而 std::stable\_sort 需要更多空间。

选择第 K 大时有三种主线。完整排序最简单，适合数据量不大或后续还需要有序结果；大小为 K 的堆适合 K 远小于 n 或数据流持续到来；快速选择适合一次性数组、需要平均线性。三种方法没有绝对优劣，差别在输出需求、数据是否在线、实现风险和最坏复杂度。高质量题解应该比较它们，而不是只写一种。

排序题的边界往往出在比较器。比较器必须满足严格弱序，不能写出自相矛盾的关系。区间覆盖题中同起点要按终点降序，合并区间按起点升序即可，会议室按开始时间或结束时间有不同含义。比较器写错时，代码看似能跑，但排序结果不满足算法不变量。每次写比较器都要用三组数据验证：相同主键、完全相同元素、主键相反顺序。

实验建议：实现完整排序、大小为 K 的堆、快速选择三种第 K 大算法，用随机数组、几乎有序数组、逆序数组和大量重复数组测试。记录比较次数或运行时间，观察 K 很小时堆的优势、需要完整有序时排序的优势、随机化快速选择的波动。再把快速选择枢轴固定为首元素，用有序数组制造最坏情况，理解为什么工程库不会使用裸快速排序。

\begin{principlebox}{本节复盘}
阅读本节后，请回到相邻章节的 LeetCode 案例，例如 LeetCode 875 爱吃香蕉、LeetCode 72 编辑距离、LeetCode 743 网络延迟时间、LeetCode 215 数组第 K 大，把暴力瓶颈、优化依据、状态不变量、C++ 容器成本和边界测试重新写一遍。能从这些具体题迁移到变体题，才算真正掌握。
\end{principlebox}
